\documentclass[12pt, a4paper]{article}
\usepackage[utf8]{inputenc}
\usepackage[T1]{fontenc}
\usepackage{geometry}
\usepackage{xcolor}
\usepackage{listings}
\usepackage{hyperref}
\usepackage{graphicx}
\usepackage{booktabs}
\usepackage{float}
\usepackage{amsmath}

% Page geometry
\geometry{top=2.5cm, bottom=2.5cm, left=2.5cm, right=2.5cm}

% Color definitions
\definecolor{codegreen}{rgb}{0,0.6,0}
\definecolor{codegray}{rgb}{0.5,0.5,0.5}
\definecolor{codepurple}{rgb}{0.58,0,0.82}
\definecolor{backcolour}{rgb}{0.95,0.95,0.92}

% Code listing style
\lstdefinestyle{mystyle}{
    backgroundcolor=\color{backcolour},   
    commentstyle=\color{codegreen},
    keywordstyle=\color{magenta},
    numberstyle=\tiny\color{codegray},
    stringstyle=\color{codepurple},
    basicstyle=\ttfamily\footnotesize,
    breakatwhitespace=false,         
    breaklines=true,                 
    captionpos=b,                    
    keepspaces=true,                 
    numbers=left,                    
    numbersep=5pt,                  
    showspaces=false,                
    showstringspaces=false,
    showtabs=false,                  
    tabsize=2
}
\lstset{style=mystyle}

\title{\textbf{System Evaluation and Measures of Success}}
\author{Development Team}
\date{}

\begin{document}

\maketitle

\begin{abstract}
This article outlines the five key measures of success used to evaluate the ``Byte Me'' food rescue platform. Each measure is defined with specific criteria and supported by evidence from the system's codebase, ensuring reliability, accuracy, and environmental impact.
\end{abstract}

\section{Marketplace Reliability}

\subsection{Oversell Rate (Target: 0\%)}
\textbf{Definition:} The system must never allow more reservations than the available quantity of a bundle.

\textbf{Evidence:} The \texttt{BundlePosting} entity enforces a strict logic check before any reservation is permitted. The \texttt{canReserve} method explicitly verifies inventory levels.

\begin{lstlisting}[language=Java, caption=BundlePosting.java Inventory Logic]
// Check if can reserve
public boolean canReserve(int qty) {
    return status == Status.ACTIVE && (quantityReserved + qty) <= quantityTotal;
}
\end{lstlisting}

Furthermore, the \texttt{OrderController} transactionally locks this process. If \texttt{canReserve(1)} returns false, the system immediately rejects the request with a \texttt{400 Bad Request}, ensuring the oversell rate remains strictly at 0.

\subsection{Handling Expiries and No-Shows}
\textbf{Definition:} The lifecycle of every bundle must be accounted for, distinguishing between successful collections, expired goods, and user no-shows.

\textbf{Evidence:} 
\begin{itemize}
    \item \textbf{Status Tracking:} The \texttt{Reservation.Status} enum includes specific states for \texttt{NO\_SHOW}, \texttt{EXPIRED}, and \texttt{CANCELLED}.
    \item \textbf{Automated Cleanup:} The \texttt{BundlePostingRepository} includes a query to identify expired bundles automatically:
\end{itemize}

\begin{lstlisting}[language=Java, caption=Finding Expired Bundles]
@Query("SELECT b FROM BundlePosting b WHERE b.status = 'ACTIVE' AND b.pickupEndAt < :now")
List<BundlePosting> findExpired(Instant now);
\end{lstlisting}

This ensures that inventory lists remain fresh and accurate without manual intervention.

\section{Forecast Accuracy}

\subsection{Methodology}
We compare our chosen predictive model against standard baselines using a Held-Out Evaluation approach (80\% training data, 20\% evaluation data).

\subsection{Metrics \& Baselines}
The \texttt{ForecastService} calculates three key error metrics for every model run:
\begin{enumerate}
    \item \textbf{MAE (Mean Absolute Error):} Average magnitude of reservation errors.
    \item \textbf{RMSE (Root Mean Squared Error):} Penalizes larger outliers.
    \item \textbf{Brier Score:} Evaluates the accuracy of probabilistic no-show predictions.
\end{enumerate}

These metrics are calculated by comparing the \textbf{Chosen Model (EMA-Weighted)} against two baselines:
\begin{itemize}
    \item \textbf{Moving Average (4-Week):} A simple sliding window average.
    \item \textbf{Seasonal Naive:} Predicts based on the same ISO Day of Week from history.
\end{itemize}

\begin{lstlisting}[language=Java, caption=Model Comparison Logic]
// Compare models and save metrics
Map<String, Object> maMetrics = Map.of(
    "MAE_reservations", calculateMAE(maPredicted, actualRes),
    "RMSE_reservations", calculateRMSE(maPredicted, actualRes),
    "Brier_no_show", calculateBrierScore(maNoShowPred, actualNoShow));
\end{lstlisting}

The system stores these results in the \texttt{ForecastRun} entity to allow for continuous performance monitoring.

\section{Waste Impact \& Carbon Reduction}

\subsection{Estimated Waste Avoided}
\textbf{Measure:} The total weight of food diverted from landfills.

\textbf{Evidence:} The \texttt{SellerMetricsWeekly} entity tracks \texttt{wasteAvoidedGrams}. This is calculated by multiplying the number of \texttt{COLLECTED} reservations by the \texttt{estimatedWeightGrams} defined in the bundle posting.

\subsection{Real-World Example: CO2e Prevention}
To illustrate, consider the sample data for ``Sourdough \& Co Bakery'' used in our forecast demo.

\begin{itemize}
    \item \textbf{Scenario:} A bakery posts surplus bundles with an estimated weight of 800g each.
    \item \textbf{Observations:} On Jan 30, 2026, the bakery observed 16 reservations with a 9\% no-show rate.
    \item \textbf{Food Waste Avoided:}
    \[
    \text{Mass} = 16 \times (1 - 0.09) \times 0.8\text{ kg} \approx 11.65\text{ kg}
    \]
    \item \textbf{Carbon Impact:} Using a standard conversion factor of 2.5 kg CO2e per kg of food waste prevented (WRAP, 2023):
    \[
    \text{CO2e Saved} = 11.65\text{ kg} \times 2.5 = \textbf{29.1 kg CO2e}
    \]
\end{itemize}

This metric is aggregated weekly and displayed on the seller dashboard, providing tangible feedback on their environmental contribution. The \texttt{CO2\_SAVER} badge is explicitly awarded when organizations reach 10 collections, further reinforcing this goal.

\section{User Experience (Task-Based)}

\subsection{Reserving and Claiming}
\textbf{Target:} Sub-second response times for reservation and collection tasks.

\textbf{Evidence:} 
\begin{itemize}
    \item \textbf{Instant Reservation:} The \texttt{create} order endpoint generates a secure 6-digit \texttt{claimCode} immediately upon success.
    \item \textbf{Secure Collection:} The collection process uses a hashed verification system. When an organisation presents a code, the backend verifies it using \texttt{passwordEncoder.matches()}.
\end{itemize}

\begin{lstlisting}[language=Java, caption=Claim Code Verification]
if (!passwordEncoder.matches(claimReq.getClaimCode(), reservation.getClaimCodeHash())) {
    return ResponseEntity.badRequest().body("Invalid claim code");
}
\end{lstlisting}

This approach balances high security with the speed required for retail environments.

\section{Gamification \& Engagement}

\subsection{Measures}
\begin{itemize}
    \item \textbf{Retention:} Tracking \texttt{currentStreakWeeks} via the \texttt{OrganisationStreakCache}.
    \item \textbf{Milestones:} Awarding badges like \texttt{STREAK\_4} and \texttt{CO2\_SAVER}.
\end{itemize}

\textbf{Evidence:} The system automatically updates streak counters and checks for badge eligibility after every successful collection, providing immediate positive reinforcement to the user.

\end{document}